\documentclass[spanish]{article}
\usepackage[utf8]{inputenc}                   % Para escribir tildes y eñes
\usepackage{fancyhdr}
\usepackage{url}
\usepackage{listings}
\usepackage{tabularx}
\usepackage{booktabs}
%\usepackage{hyperref}
\usepackage[colorlinks=true,urlcolor=blue,linkcolor=blue,citecolor=blue]{hyperref}
\usepackage{graphicx}  
\usepackage{pgfgantt}
\hyphenation{}
\usepackage[vmargin=4cm,			%%margins
	    tmargin=3cm,
	    hmargin=2cm,
	    letterpaper]{geometry}

\newcommand{\entitle}[1]{
  \vspace{0.3cm}%
  \noindent%
  \textbf{#1}%
  \vspace{0.3cm}%
  \hrule\vspace*{0.5mm}%
  \noindent%
  \rule{\linewidth}{0.5mm}%
  \vspace{0.5cm}%
}

\begin{document}
\fancyhead[L]{\includegraphics[scale=0.18]{./images/ucr}}
\fancyhead[C]{UNIVERSIDAD DE COSTA RICA\\ESCUELA DE INGENIERÍA ELECTRICA\\ESTRUCTURAS ABSTRACTAS DE DATOS Y ALGORITMOS\\ PARA INGENIERÍA\\PROPUESTA LIBRERÍA CRIPTOGRÁFICA EN C++\\GEYKE93
}
\fancyhead[R]{\includegraphics[scale=0.45]{./images/eie}}
\pagestyle{fancy}
\setlength{\headheight}{60pt}

\entitle{
\begin{center}
GABRIEL JIMÉNEZ LEIVA B23466 \\
CARLA VEGA JARQUÍN B06763\\
\end{center}
}

\section{Justificación}
La aparición de la informática y el uso masivo de las comunicaciones digitales han producido un número creciente de problemas de seguridad. Las transacciones que se realizan a través de la red pueden ser interceptadas. La seguridad de esta información debe garantizarse. Este desafío lleva al estudio de los algoritmos, protocolos y sistemas que se utilizan para proteger la información y dotar de seguridad a las comunicaciones y a las entidades que se comunican. Los criptógrafos investigan, desarrollan y aprovechan técnicas matemáticas que les sirven como herramientas para conseguir sus objetivos. 

\section{Metodología}
En este trabajo se busca partir del método de encriptación basado en los números primos y ejemplificar sus implementaciones en la criptografía antigua y la moderna. Dentro de esta última, estudiaremos a Ronald Linn Rivest y NSA (National Security Agency) que han sido pioneros de la época y quienes cuyo trabajo ha cambiado y clasificado la criptografía a través métodos que pensamos utilizar para la creación de la librería que es el fin del proyecto.

\subsection{Tabla de algoritmos}
\begin{table}[h]
\begin{center}
\begin{tabular}{c | c | c}
\toprule[0.3mm]
Encriptacion simétrica & Encriptacion asimétrica & Hash function \\
\midrule
AES (Rijndael) & RSA & SHA2(256-512)\\
\bottomrule[0.3mm]
\end{tabular}
\end{center}
\caption{Métodos de encriptación}
\end{table}


\section{Objetivos}

\subsection{Objetivo General}
Crear una librería de encriptacion que proporcione métodos de encriptación simétrica (AES), asimétrica (RSA) y funciones hash (SHA2). 

\subsection{Objetivos Específicos}
\begin{itemize}
\item Implementar un algoritmo de encriptación y desencriptación simétrica AES (Rijndael) y otro asimétrica RSA utilizando conocimientos en el lenguaje de programacion C++: operaciones booleanas, funciones algebraicas, cifrado por bloques y otras herramientas útiles para la criptografía.
\item Implementar una función hash SHA2, tanto en su versión SHA-256 para sistemas de 32 bits como SHA-512 para 64 bits.
\item Mencionar los algoritmos criptográficos y funciones hash de librerías como Crypto++, Beecrypt y hashlib++.
\item Estudiar las implementaciones de algoritmos criptográficos y funciones hash en un sistema como SSH.
\end{itemize}

\section{Referencias}
\begin{thebibliography}{10}

\bibitem{}\url{http://beecrypt.sourceforge.net/doxygen/c++/index.html}
\bibitem{}\url{http://sourceforge.net/projects/cryptopp/?source=recommended}
\bibitem{}\url{http://sourceforge.net/projects/hashlib2plus/?source=recommended}
\bibitem{}\url{http://www.aescrypt.com/}
\bibitem{}\url{http://es.wikipedia.org/wiki/Ronald_Rivest}
\bibitem{}\url{http://es.wikipedia.org/wiki/SSH}
\bibitem{}\url{http://es.wikipedia.org/wiki/RSA}
\bibitem{}\url{http://es.wikipedia.org/wiki/Funci%C3%B3n_hash}
\bibitem{}\url{http://www.openssh.org/}
\bibitem{}\url{mirror.math.ku.edu/tex.../pgfgantt/pgfgantt.pdf}

\section{Anexo}

La distribución del trabajo a través del diagrama de Gantt se planificó de la siguiente manera: \\

\begin{ganttchart}{1}{11}
\gantttitle{Setiembre y Octubre}{11} \\
\gantttitlelist{23,24,25,26,27,28,29,30,1,2,3}{1} \\
\ganttgroup{Hash function (Gabriel)}{1}{7} \\
\ganttgroup{AES (Carla)}{1}{7}\\
\ganttgroup{RSA (Gabriel y Carla)}{8}{11} \\
\end{ganttchart}

\end{thebibliography}

\end{document}


